\documentclass[11pt]{article}

% ============================================================
% AI Co-Mathematician — working paper
% Project: {{PROJECT_NAME}}
% Initialized: {{DATE}}
% ============================================================

\usepackage[a4paper,margin=1in,marginparwidth=1.4in,marginparsep=0.2in]{geometry}
\usepackage{amsmath,amssymb,amsthm}
\usepackage{mathtools}
\usepackage{hyperref}
\usepackage{xcolor}
\usepackage{marginnote}
\usepackage{enumitem}
\usepackage{etoolbox}

\hypersetup{
  colorlinks=true,
  linkcolor=blue!50!black,
  citecolor=blue!50!black,
  urlcolor=blue!50!black
}

% ---- Theorem environments --------------------------------------------------
\newtheorem{theorem}{Theorem}[section]
\newtheorem{lemma}[theorem]{Lemma}
\newtheorem{proposition}[theorem]{Proposition}
\newtheorem{corollary}[theorem]{Corollary}
\theoremstyle{definition}
\newtheorem{definition}[theorem]{Definition}
\newtheorem{example}[theorem]{Example}
\theoremstyle{remark}
\newtheorem{remark}[theorem]{Remark}

% ---- Co-mathematician custom macros ---------------------------------------
%
% \unproven{text}     — flags an unverified step. Renders in red and is
%                       collected in the "Open Obligations" appendix.
% \marginorigin{tag}  — provenance tag in the margin (e.g., workstream id,
%                       sub-agent name, citation source).
% \internalref{path}  — link to a file in the project (workstream report,
%                       failed exploration, etc.).
% \claim{text}        — an asserted claim that needs justification or citation.
%                       Reviewer agent must check that every \claim is
%                       followed by a proof, citation, or \unproven.

\newcounter{unprovencount}
\newcommand{\unprovenlist}{}

\newcommand{\unproven}[1]{%
  \stepcounter{unprovencount}%
  \textcolor{red!80!black}{\textbf{[UNPROVEN \theunprovencount]} #1}%
  \marginnote{\footnotesize\textcolor{red!80!black}{unproven \#\theunprovencount}}%
  \xappto\unprovenlist{\item[\textbf{\#\theunprovencount}] #1\par}%
}

\newcommand{\marginorigin}[1]{\marginnote{\footnotesize\textit{origin:} #1}}

\newcommand{\internalref}[2]{\href{run:#1}{\texttt{#2}}}

\newcommand{\claim}[1]{\textit{#1}\marginnote{\footnotesize\textcolor{orange!70!black}{claim — needs justification}}}

% \leanproved{workstream-id} — placed inside (or right after) a theorem
% environment in lieu of \proof, indicating the result has been verified by
% the lean-prover sub-agent. The hook paper_tex_guard.py treats this as a
% valid alternative to \proof. Links the reader to the Lean source.
\newcommand{\leanproved}[1]{%
  \textcolor{green!50!black}{$\blacksquare$ \textit{Lean-verified in }\internalref{workstreams/#1/lean/Main.lean}{\texttt{#1}}\textit{.}}%
  \marginnote{\footnotesize\textcolor{green!50!black}{Lean: #1}}%
}

% ---- Document --------------------------------------------------------------
\title{ {{PROJECT_NAME}} \\ \large (working paper) }
\author{ }
\date{Initialized {{DATE}}}

\begin{document}
\maketitle

\begin{abstract}
\textit{Working draft.} Research question: {{RESEARCH_QUESTION}}

\medskip
This is a living document maintained by the AI co-mathematician system. Sections
will be added, revised, and reviewed iteratively. Margin notes record the
provenance of each contribution. Items flagged \unproven{} are explicitly
unverified and listed in the appendix of open obligations.
\end{abstract}

\tableofcontents

% ---- Sections start here --------------------------------------------------

\section{Introduction}

\marginorigin{co-math-init scaffolding}
This section will be drafted by the project-coordinator agent after the user
approves the goals in \internalref{goals.md}{goals.md}.

\section{Preliminaries}

\textit{To be drafted by literature-reviewer workstream.}

\section{Main results}

\textit{To be drafted by prover workstream(s).}

\section{Computational evidence}

\textit{To be drafted by coder workstream(s).}

\section{Discussion}

\textit{To be drafted last, after main results stabilise.}

% ---- Open obligations appendix --------------------------------------------

\appendix

\section{Open obligations}
\label{app:open}

The following items are flagged \texttt{\textbackslash{}unproven} elsewhere in
this document. The reviewer agent must verify that each is either resolved or
explicitly accepted before the paper can be marked complete.

\begin{description}
\unprovenlist
\end{description}

\section{Failed explorations (negative space)}

Per the co-mathematician principle of preserving the history of unsuccessful
attempts, summaries of abandoned approaches live in
\internalref{failed-explorations/}{failed-explorations/}. Reviewer agents
consult this directory before suggesting new strategies.

\end{document}
